\noindent\textbf{ZUR NETZAUSGLEICHUNG.}\hfill323
\begin{center}\textsc{Gauss} an \textsc{Gerling}.\quad Göttingen, 2. Mai 1837.\end{center}

.\ .\ .\ .\ .\ .\ Was die Berechnung der geodätischen Messungen betrifft, so
habe ich zwar die Elimination bei meinen eigenen Hauptdreiecken von Göttingen
bis Jever, gegen 40 Dreiecke zwischen 33 Punkten, ganz vollständig und
nach aller Strenge ausgeführt. Es blieb dies aber damals eine sehr beschwerliche
langwierige Arbeit trotz der Kunstgriffe, die ich dabei freilich angewandt
habe. Allein diese Ihnen genügend zu erklären, würde ein kleines Buch, und
um dieses auszuarbeiten, erst ein Wiederhineinstudiren von meiner Seite nöthig
sein, woran ich jetzt gar nicht denken kann.

\begingroup
\def\letterfourstars{*}
\makeatletter
\renewcommand{\@makefnmark}{\hbox{\normalfont[\textsuperscript{\letterfourstars})]}}
\renewcommand{\@makefntext}[1]{\parindent1em\noindent\hbox to2em{\hss[\textsuperscript{\letterfourstars})\ }#1]}
\makeatother
Bei den später von meinem Sohn, \textsc{Hartmann} und \textsc{Müller} ausgeführten
Dreiecken, in Westphalen, Lüneburgschen, Harz und Wesergegend, sowie
den vorjährigen, habe ich von der äussersten Strenge etwas nachgelassen, und
es so gemacht, dass ich
\begin{list}{}{\setlength{\leftmargin}{2.4em}\setlength{\labelwidth}{1.5em}\setlength{\labelsep}{.5em}\setlength{\itemsep}{.3em}\setlength{\topsep}{.3em}}
\item[1)] die Winkelbedingungen \textls[120]{allein} in Betracht zog und danach scharf ausglich. Dann
\item[2)] zu diesen ausgeglichenen Zahlen die neue Ausgleichung suchte, welche
die Seitenverhältnisse erfordern\footnote{Vergl. S.\ 253.}.
\end{list}
\def\letterfourstars{**}
\par\noindent\hspace*{\parindent}Will man sich die Mühe geben, dies Verfahren alternirend wiederholt
anzuwenden, bis man zu stehenden Resultaten kommt, was ich bereits im
Supplementum etc.\ p.\ 28\footnote{Band IV, S.\ 77/79.} anempfohlen habe, so gelangt man, wie dort
bewiesen ist, \textls[120]{genau} zu denselben Resultaten, als wenn man alle Bedingungen
auf Einmal berücksichtigte.
\par\medskip
.\ .\ .\ .\ .\ .
\endgroup
